\chapter{Teoria della probabilità in spazi finiti}

La teoria matematica della probabilità mira a costruire un formalismo teorico per descrivere i \emph{fenomeni aleatori} (casuali). Si basa sul concetto di \emph{esperimento aleatorio}, che può essere approssimativamente definito come "un esperimento il cui risultato non può essere previsto in anticipo". Un individuo può essere naturalmente portato ad associare un \emph{grado di fiducia} ai diversi possibili risultati di questo esperimento. Vediamo come modellizzare matematicamente questa idea, tenendo sempre presente che in ogni buona teoria vi sono tre fasi distinte:
\begin{enumerate}[(i)]
    \item Scelta del modello matematico;
    \item Elaborazione matematica all'interno del modello scelto;
    \item Interpretazione dei risultati ottenuti.
\end{enumerate}

\section{Spazi di probabilità}

La struttura degli spazi di probabilità si compone di tre ingredienti fondamentali:
\begin{enumerate}
    \item Lo spazio campionario (o spazio degli stati) $\Omega$;
    \item L'insieme degli eventi $\F$;
    \item La misura di probabilità $\PP$.
\end{enumerate}
Di conseguenza, uno spazio di probabilità è una terna: $(\Omega, \F, \PP)$. Descriviamo nel dettaglio gli oggetti di questa terna.

\begin{hint}[La Terna Fondamentale]
    Per ricordare la terna $(\Omega, \F, \PP)$:
    \begin{itemize}
        \item $\Omega$ (Omega) è \textbf{TUTTO} quello che può succedere (lo Spazio).
        \item $\F$ (F corsiva) sono le \textbf{DOMANDE} che puoi fare (gli Eventi a cui puoi assegnare una probabilità).
        \item $\PP$ (P) è la \textbf{RISPOSTA} (la Misura, il numero da 0 a 1).
    \end{itemize}
\end{hint}

\subsection{Spazio campionario ($\Omega$)}

Lo \emph{spazio campionario} è l'insieme i cui elementi rappresentano tutti i possibili "stati futuri/sconosciuti del mondo" \emph{mutuamente esclusivi}. È tradizionalmente indicato con $\Omega$ e il suo elemento generico è indicato con $\omega$. In questo capitolo assumiamo che lo spazio campionario $\Omega$ sia finito, ovvero $|\Omega| < \infty$. Naturalmente questa è un'assunzione restrittiva, ma ci permette comunque di modellare molti fenomeni.

\begin{esempio}[Monete e Dadi]
    \begin{itemize}
        \item \textbf{Lancio di una moneta:} Una scelta ragionevole per lo spazio campionario è $\Omega = \{T, C\}$ (Testa, Croce) oppure $\Omega = \{0, 1\}$. Le etichette non sono importanti.
        \item \textbf{Due lanci di una moneta (o lancio di due monete):} Una scelta logica è elencare le coppie: $\Omega := \{(T, T), (T, C), (C, T), (C, C)\}$.
        \item \textbf{Lancio di due dadi a sei facce:} $\Omega := \{(i, j) : i, j \in \{1, 2, 3, 4, 5, 6\}\}$. Si noti che $|\Omega| = 36$.
    \end{itemize}
\end{esempio}

\subsection{Spazio degli eventi ($\F$)}

Matematicamente, un \emph{evento} è un sottoinsieme $A \subseteq \Omega$ di risultati $\omega$. Interpretiamo $A$ come un evento che \emph{si è verificato} se il risultato effettivo $\omega \in A$, e che \emph{non si è verificato} se $\omega \notin A$. Questa rappresentazione matematica ci permette di tradurre domande intuitive in operazioni insiemistiche precise:
\begin{itemize}
    \item L'evento $\Omega$ è l'evento \emph{certo}.
    \item L'evento $\emptyset$ è l'evento \emph{impossibile}.
    \item Il singoletto $\{\omega\}$ è un evento \emph{elementare}.
    \item Dato $A \subseteq \Omega$, $A^c$ è l'"evento complementare (o contrario) di $A$".
    \item L'unione $A \cup B$ è l'evento "$A$ e/o $B$ si verificano".
    \item L'intersezione $A \cap B$ è l'evento "Sia $A$ che $B$ si verificano contemporaneamente".
\end{itemize}

L'insieme di tutti gli eventi di interesse è solitamente denotato con $\F$. Questo insieme deve avere delle "proprietà di stabilità", ovvero deve essere un'\emph{algebra}.

\begin{definition}[Algebra di insiemi e atomi]
Un insieme di sottoinsiemi $\F \subseteq 2^\Omega$ è chiamato \emph{algebra (di insiemi)} se valgono le seguenti proprietà:
\begin{enumerate}[(F1)]
    \item $\Omega \in \F$;
    \item $A \in \F \implies A^c \in \F$;
    \item $A, B \in \F \implies A \cup B \in \F$.
\end{enumerate}
Se $\F$ è un'algebra, un insieme $A \in \F$ è chiamato \emph{atomo} di $\F$ se non può essere decomposto come unione di due insiemi non vuoti e disgiunti di $\F$.
\end{definition}

\begin{spiegazione}
    Perché ci serve un'algebra? Perché se possiamo chiederci qual è la probabilità che piova (evento $A$), dobbiamo poterci chiedere anche qual è la probabilità che NON piova (evento $A^c$), oppure che piova o nevichi ($A \cup B$). L'algebra garantisce che il nostro "vocabolario" di eventi sia logicamente completo.
\end{spiegazione}

\subsection{Misura di probabilità ($\PP$)}

Una volta fissato lo spazio degli eventi $\F$, ad ogni evento $A \in \F$ si può associare un "grado di fiducia" chiamato \emph{probabilità dell'evento A}. Questo è formalizzato da una funzione:
\[ \PP : \F \to [0, 1] \]
La scelta di questa misura è soggettiva, ma deve rispettare i "tre assiomi della probabilità" per escludere scelte illogiche.

\begin{definition}[Misura di probabilità]
Sia $\Omega$ un insieme finito e sia $\F \subseteq 2^\Omega$ un'algebra. Una \emph{misura di probabilità} su $\F$ è una mappa $\PP : \F \to \R$ tale che:
\begin{enumerate}[(A1)]
    \item $0 \le \PP(A) \le 1$ per ogni $A \in \F$;
    \item $\PP(\Omega) = 1$;
    \item $A, B \in \F$ con $A \cap B = \emptyset \implies \PP(A \cup B) = \PP(A) + \PP(B)$.
\end{enumerate}
\end{definition}

\begin{hint}[Additività]
    L'assioma (A3) è la \textbf{Regola della Somma}: se due eventi non possono verificarsi insieme (sono disgiunti), la probabilità che se ne verifichi almeno uno è la somma delle loro probabilità singole. (Es. Probabilità di tirare 1 o 2 con un dado = P(1) + P(2)).
\end{hint}

Dagli assiomi derivano proprietà utilissime (prova a dimostrarle mentalmente usando i diagrammi di Eulero-Venn):
\begin{itemize}
    \item $\PP(A^c) = 1 - \PP(A)$
    \item Se $A \subseteq B$, allora $\PP(B \setminus A) = \PP(B) - \PP(A)$ e quindi $\PP(A) \le \PP(B)$.
    \item $\PP(A \cup B) = \PP(A) + \PP(B) - \PP(A \cap B)$
\end{itemize}

\section{Variabili aleatorie}

Nelle applicazioni pratiche, spesso non ci interessa l'esito esatto $\omega$, ma un numero o una quantità derivata da esso (es. la somma di due dadi, il guadagno in un gioco). Questo ci porta alle variabili aleatorie.

\begin{definition}[Variabile aleatoria]
Sia dato uno spazio $(\Omega, \F)$. Una mappa $X : \Omega \to E$ è una \emph{variabile aleatoria} se è $\F$-misurabile, ovvero se:
\[ X^{-1}(B) := \{X \in B\} \in \F \quad \forall B \subseteq E \]
\end{definition}

In spazi campionari finiti in cui $\F = 2^\Omega$ (l'insieme delle parti, ovvero tutti i possibili sottoinsiemi), \textbf{ogni funzione $X$ è una variabile aleatoria}.

\begin{spiegazione}
    Una variabile aleatoria non è né una "variabile" in senso stretto, né tantomeno "aleatoria" (casuale). È semplicemente una \textbf{funzione deterministica} che prende in input lo stato del mondo $\omega$ (quello sì che è casuale) e restituisce un valore, di solito un numero. \\
    Esempio: se $\omega = (\text{Testa}, \text{Testa})$, la variabile aleatoria $X = \text{"numero di teste"}$ restituirà banalmente $X(\omega) = 2$.
\end{spiegazione}

\subsection{Misurabilità di una mappa e algebra generata}

Il concetto di "algebra generata" serve a formalizzare l'idea di "informazione trasportata da una variabile aleatoria".

\begin{definition}[Algebra generata $\sigma(X)$]
Sia $X : \Omega \to E$ una variabile aleatoria. L'\emph{algebra generata} da $X$, denotata con $\sigma(X)$, è la più piccola algebra su $\Omega$ che rende $X$ misurabile:
\[ \sigma(X) := \{ X^{-1}(B) : B \subseteq E \} \]
\end{definition}

Gli atomi di $\sigma(X)$ sono esattamente gli insiemi di livello $\{X = x\}$ per ogni $x \in E$. 
\textbf{Intuizione:} L'algebra $\sigma(X)$ rappresenta tutta l'informazione che deduco osservando $X$.

\begin{theorem}[Criterio di misurabilità di Doob]
Siano $X : \Omega \to E$ e $Y : \Omega \to E'$ variabili aleatorie. 
$X$ è $\sigma(Y)$-misurabile se e solo se esiste una funzione $f : E' \to E$ tale che $X = f \circ Y$.
\end{theorem}

\begin{hint}[Interpretare Doob]
    Cosa significa che $X$ è $\sigma(Y)$-misurabile? 
    Significa che \textbf{l'informazione fornita da $Y$ è sufficiente a determinare $X$}. Se conosco $Y$, conosco automaticamente anche $X$. 
    Perché? Perché se $X = f(Y)$, basta applicare la funzione deterministica $f$ al valore di $Y$ per trovare $X$. Se invece serve lanciare una moneta aggiuntiva, allora $X$ \textbf{non} è $\sigma(Y)$-misurabile.
\end{hint}

\subsection{Legge (o distribuzione) di una variabile aleatoria}

\begin{definition}[Legge o distribuzione]
Data una variabile aleatoria $X : \Omega \to E$, possiamo trasferire la misura di probabilità da $(\Omega, \F, \PP)$ allo spazio dei valori $(E, 2^E)$ definendo:
\[ \mu_X(B) := \PP(X \in B) \quad \forall B \subseteq E \]
Questa nuova misura di probabilità si chiama \emph{legge} o \emph{distribuzione} di $X$.
La funzione $f_X(x) := \PP(X = x)$ è la sua \emph{funzione di densità discreta}.
\end{definition}

Poiché stiamo assumendo che $E$ sia finito, la distribuzione di probabilità (o legge) di una variabile aleatoria $X : \Omega \to E$ è completamente determinata dalla sua funzione di densità discreta $f_X$. Infatti, per ogni $B \subseteq E$, possiamo recuperare la legge $\mu_X(B)$ semplicemente sommando:
\[ \mu_X(B) = \sum_{x \in B} f_X(x) \]
Questa proprietà additiva rende le distribuzioni su spazi finiti particolarmente trattabili. Mentre la legge fornisce la misura di probabilità astratta, la funzione di densità discreta ci dà uno strumento computazionale concreto.

\begin{esempio}[Somma di due dadi]
Riprendendo il lancio di due dadi a sei facce, consideriamo la variabile aleatoria $S(\omega) = i + j$ (la somma dei due dadi). Il codominio è $E = \{2, \dots, 12\}$. La legge $\mu_S$ è descritta dalla densità discreta $f_S$ calcolata contando i casi favorevoli:
\[ f_S(2) = \frac{1}{36}, \quad f_S(3) = \frac{2}{36}, \quad \dots \quad f_S(7) = \frac{6}{36}, \dots \]
\end{esempio}

\subsection{Distribuzione congiunta e distribuzioni marginali}

Quando studiamo più variabili aleatorie contemporaneamente, abbiamo bisogno di capire non solo il loro comportamento individuale, ma anche come si relazionano tra loro.

\begin{definition}[Leggi congiunte e marginali]
Date $n$ variabili aleatorie $X_j : \Omega \to E_j$, con $j=1,\dots,n$, il vettore $X = (X_1, \dots, X_n)$ è una variabile aleatoria che assume valori nello spazio prodotto $E = E_1 \times \dots \times E_n$.
La \emph{legge congiunta} $\mu_X = L(X_1, \dots, X_n)$ è la misura di probabilità su $E$ definita da:
\[ \mu_X(A_1 \times \dots \times A_n) := \PP(X_1 \in A_1, \dots, X_n \in A_n) \]
Le \emph{leggi marginali} $\mu_{X_i} = L(X_i)$ si ottengono dalla legge congiunta:
\[ \mu_{X_i}(A_i) = \mu_X(E_1 \times \dots \times A_i \times \dots \times E_n) \]
La corrispondente \emph{densità discreta marginale} di $X_i$ si ottiene sommando rispetto a tutte le altre variabili:
\[ f_{X_i}(x_i) = \sum_{x_1, \dots, x_{i-1}, x_{i+1}, \dots, x_n} f_X(x_1, \dots, x_n) \]
\end{definition}

\begin{spiegazione}
    Immagina la legge congiunta come una grande tabella a doppia entrata (se abbiamo due variabili $X$ e $Y$). Le celle interne della tabella contengono la probabilità che $X=x$ e $Y=y$ contemporaneamente. 
    Se sommiamo tutti i valori di una riga, troviamo la probabilità totale di $X=x$ (a prescindere da $Y$). Questa somma "ai margini" della tabella è proprio la \emph{distribuzione marginale}.
\end{spiegazione}

\begin{hint}[Congiunta vs Marginale]
    Ricorda: \textbf{Dalla congiunta puoi SEMPRE ricavare le marginali (sommando), ma dalle marginali NON puoi ricavare la congiunta} (a meno che le variabili non siano indipendenti). Diverse distribuzioni congiunte possono avere le stesse identiche distribuzioni marginali.
\end{hint}

\section{Probabilità Condizionata}

Esploriamo ora due concetti fondamentali che catturano come le probabilità cambiano quando arrivano nuove informazioni.

\begin{definition}[Probabilità Condizionata]
Dato uno spazio di probabilità $(\Omega, \F, \PP)$ e un evento $H \in \F$ con $\PP(H) > 0$, la \emph{probabilità condizionata} di $A$ dato $H$ è definita come:
\[ \PP(A|H) := \frac{\PP(A \cap H)}{\PP(H)} \]
\end{definition}

La probabilità condizionata rappresenta l'aggiornamento razionale delle credenze di un agente quando viene a sapere che l'evento $H$ si è verificato. L'operazione "rinormalizza" la probabilità originale concentrandola su $H$.

\begin{esempio}[Il problema di Monty Hall]
Sei in un gioco a premi e devi scegliere tra 3 porte: dietro una c'è un'auto, dietro le altre due ci sono capre. Scegli la porta 1. Il presentatore (che sa cosa c'è dietro le porte) apre la porta 3, mostrando una capra. Ti chiede: "Vuoi cambiare e scegliere la porta 2?".
\textbf{Conviene cambiare?}
Sì, le probabilità raddoppiano (da 1/3 a 2/3). Il motivo è che l'azione del presentatore non è casuale: è costretto ad aprire una porta con una capra. Se la tua prima scelta era sbagliata (probabilità 2/3), il presentatore ti sta indicando esattamente l'altra porta sbagliata, lasciando l'auto dietro l'unica porta chiusa rimasta. Cambiando vinci sempre in questo caso!
\end{esempio}

\subsection{Legge della probabilità totale e Teorema di Bayes}

Queste formule sono fondamentali per scomporre eventi complessi in parti più semplici o per invertire i condizionamenti.

\begin{theorem}[Probabilità totale e Bayes]
\begin{enumerate}
    \item \textbf{Legge della probabilità totale}: Se $A_1, \dots, A_n$ formano una partizione di $\Omega$, per ogni evento $A$:
    \[ \PP(A) = \sum_{i=1}^n \PP(A_i) \PP(A|A_i) \]
    \item \textbf{Teorema di Bayes}: Per due eventi $A$ e $B$ con $\PP(A)>0$ e $\PP(B)>0$:
    \[ \PP(A|B) = \frac{\PP(B|A)\PP(A)}{\PP(B)} \]
\end{enumerate}
\end{theorem}

\begin{spiegazione}
    \textbf{Probabilità Totale}: Se voglio calcolare la probabilità che piova (evento A), posso dividere il mondo in due scenari (partizione): $A_1$=è nuvoloso, $A_2$=è sereno. La probabilità che piova sarà la probabilità che piova QUANDO è nuvoloso (moltiplicata per la probabilità che sia nuvoloso) PIÙ la probabilità che piova QUANDO è sereno (moltiplicata per la probabilità che sia sereno). \\
    \textbf{Bayes}: Serve a "ribaltare" la domanda. Conosciamo la probabilità di avere un test positivo sapendo di essere malati $P(+|Malato)$, ma a noi interessa la probabilità di essere malati avendo il test positivo $P(Malato|+)$! Bayes fa esattamente questo.
\end{spiegazione}

\begin{hint}[Formula di Bayes in Pratica]
    Quando applichi Bayes, spesso il denominatore $\PP(B)$ non è dato direttamente. Ricorda di espanderlo sempre usando la Legge della Probabilità Totale al denominatore: 
    \[ \PP(A_1|B) = \frac{\PP(B|A_1)\PP(A_1)}{\PP(B|A_1)\PP(A_1) + \PP(B|A_2)\PP(A_2) + \dots} \]
\end{hint}

\section{Indipendenza}

L'indipendenza cattura le situazioni in cui la conoscenza di un fenomeno casuale non fornisce alcuna informazione su un altro. Intuitivamente, non c'è "interazione".

\begin{definition}[Indipendenza tra variabili aleatorie]
Due variabili aleatorie $X : \Omega \to E$ e $Y : \Omega \to E'$ sono \emph{indipendenti} se per ogni evento della forma $H = \{Y \in B\}$ con $B \subseteq E'$ e $\PP(H)>0$, si ha:
\[ \PP(X \in A | Y \in B) = \PP(X \in A) \]
In modo equivalente, c'è una fattorizzazione simmetrica:
\[ \PP(X \in A, Y \in B) = \PP(X \in A)\PP(Y \in B) \]
\end{definition}

\begin{definition}[Indipendenza tra eventi]
Due eventi $A, B \in \F$ sono indipendenti se e solo se:
\[ \PP(A \cap B) = \PP(A)\PP(B) \]
Questo equivale a dire che le loro funzioni indicatrici $\mathbf{1}_A$ e $\mathbf{1}_B$ sono variabili aleatorie indipendenti.
\end{definition}
\begin{theorem}[Caratterizzazione dell'indipendenza]
Due variabili aleatorie $X$ e $Y$ sono indipendenti se e solo se la loro distribuzione congiunta è il prodotto delle marginali:
\[ \mu_{(X,Y)}(A \times B) = \mu_X(A)\mu_Y(B) \]
Oppure, in termini di densità:
\[ f_{X,Y}(x, y) = f_X(x)f_Y(y) \]
\end{theorem}

\subsection{Famiglie indipendenti}

L'indipendenza a due a due non basta quando abbiamo molte variabili.

\begin{definition}[Indipendenza di una famiglia]
Una famiglia finita di variabili aleatorie $\{X_1, \dots, X_n\}$ è \emph{indipendente} se per qualsiasi due sottofamiglie disgiunte, i vettori risultanti sono indipendenti. Questo equivale a dire che la densità congiunta fattorizza in tutte le marginali:
\[ f_{X_1, \dots, X_n}(x_1, \dots, x_n) = f_{X_1}(x_1) \dots f_{X_n}(x_n) \]
\end{definition}

\begin{spiegazione}
    Un errore comune è pensare che se tutte le coppie di variabili in un gruppo sono indipendenti tra loro (indipendenza a due a due), allora tutto il gruppo è indipendente. L'esempio di Bernstein ci mostra che questo è falso! Due monete possono essere indipendenti, ma una terza moneta il cui risultato dipende dal prodotto delle prime due rompe l'indipendenza globale, pur mantenendola a coppie.
\end{spiegazione}

\begin{theorem}[Conservazione dell'indipendenza per trasformazioni]
Se $\{X_1, \dots, X_n\}$ sono variabili indipendenti, allora applicando delle funzioni deterministiche (es. elevarle al quadrato, moltiplicarle per costanti), le nuove variabili risultanti $\{f_1(X_1), \dots, f_n(X_n)\}$ \textbf{rimangono indipendenti}. L'elaborazione deterministica non "crea" correlazione dal nulla.
\end{theorem}

\section{Tre distribuzioni fondamentali (su insiemi finiti)}

\subsection{Distribuzione Uniforme}
Modella esperimenti in cui un oggetto è scelto in modo del tutto casuale ed equo da un insieme di $n$ elementi.
La densità discreta è:
\[ f_X(x) = \frac{1}{n} \]

\subsection{Distribuzione di Bernoulli}
Modella un singolo esperimento binario (es. lancio di moneta), con esito 1 (successo) o 0 (fallimento).
Indicata con $\mathcal{B}(p)$ dove $p \in (0,1)$ è la probabilità di successo.
\[ f_X(1) = p, \quad f_X(0) = 1-p \]

\subsection{Distribuzione Binomiale}
Conta il \textbf{numero totale di successi} su $n$ esperimenti di Bernoulli \emph{indipendenti} e con la stessa probabilità $p$.
Denotata con $\mathcal{B}(n, p)$.
Per calcolare la probabilità di avere esattamente $i$ successi, si usa il calcolo combinatorio (in particolare il coefficiente binomiale $\binom{n}{i}$):
\[ \PP(S = i) = \binom{n}{i} p^i (1-p)^{n-i} \]

\begin{hint}[Binomiale = Somma di Bernoulli]
    È fondamentale ricordare che una variabile Binomiale $S \sim \mathcal{B}(n,p)$ può SEMPRE essere scritta come somma di $n$ variabili di Bernoulli indipendenti: $S = X_1 + X_2 + \dots + X_n$. Questo "trucco" rende i calcoli (es. il valore atteso) facilissimi usando la linearità!
\end{hint}

\section{Valore Atteso e Varianza}

Il \emph{valore atteso} (o media) è il concetto fondamentale per quantificare la "tendenza" centrale di una variabile aleatoria.

\begin{definition}[Valore atteso]
Sia $X : \Omega \to \R$ una variabile aleatoria. Il suo valore atteso è:
\[ \E[X] := \sum_{x \in E} x \PP(X = x) = \sum_{x \in E} x f_X(x) \]
\end{definition}

\begin{spiegazione}
    Il valore atteso NON è il valore che ti "aspetti" di vedere in un singolo esperimento (può anche essere un numero impossibile da ottenere, come 3.5 per un dado). È la \textbf{media ponderata} di tutti i possibili valori, dove i pesi sono le probabilità. Rappresenta la migliore stima costante della variabile.
\end{spiegazione}

Le proprietà chiave del valore atteso sono la \textbf{linearità} e la \textbf{monotonia}:
\[ \E[\alpha X + \beta Y] = \alpha \E[X] + \beta \E[Y] \]

\begin{definition}[Varianza]
La varianza misura quanto una variabile si "discosta" dal suo valore medio. È definita come:
\[ \operatorname{Var}(X) := \E[(X - \E[X])^2] \]
\end{definition}

\section{Valore Atteso Condizionato}

Consideriamo ora il valore atteso di una variabile aleatoria rispetto a una probabilità condizionata.

\begin{definition}[Valore atteso condizionato a un evento]
Dato un evento $H \in \F$ con $\PP(H) > 0$, il \emph{valore atteso condizionato} di $X$ dato $H$ è:
\[ \E[X|H] := \sum_{x \in E} x \PP(X = x | H) \]
\end{definition}

Questo ci permette di definire una formula utilissima, speculare alla Legge della Probabilità Totale.

\begin{theorem}[Formula di disintegrazione]
Se $A_1, \dots, A_n$ è una partizione di $\Omega$, il valore atteso globale si può ricavare combinando i valori attesi condizionati:
\[ \E[X] = \sum_{i=1}^n \E[X|A_i] \PP(A_i) \]
\end{theorem}

\subsection{Valore atteso condizionato rispetto a una variabile aleatoria}

Possiamo anche condizionare rispetto a un'altra variabile aleatoria $Y$. In questo caso, il risultato non è un singolo numero, ma \textbf{una nuova variabile aleatoria} che dipende dai valori assunti da $Y$.

\begin{definition}
Il \emph{valore atteso condizionato} di $X$ data $Y$, denotato con $\E[X|Y]$, è la variabile aleatoria definita come:
\[ \E[X|Y](\omega) = \E[X | Y = Y(\omega)] \]
\end{definition}

\begin{spiegazione}
    Mentre $\E[X|Y=y]$ è un numero fisso (perché abbiamo "fissato" l'osservazione $Y=y$), $\E[X|Y]$ è una funzione deterministica di $Y$. Rappresenta la nostra "migliore stima" di $X$ in base a quale valore assume $Y$ nel mondo.
\end{spiegazione}

Proprietà fondamentali di $\E[X|Y]$:
\begin{itemize}
    \item Se $X$ e $Y$ sono indipendenti: la conoscenza di $Y$ non aiuta a stimare $X$, quindi $\E[X|Y] = \E[X]$.
    \item Se $X$ è completamente determinata da $Y$ (cioè $X = f(Y)$): la stima migliore è il valore esatto, quindi $\E[X|Y] = X$.
    \item \textbf{Legge delle aspettative iterate (Tower property)}: $\E[\E[X|Y]] = \E[X]$. In media, la nostra stima condizionata ci restituisce il valore atteso globale!
\end{itemize}

\begin{theorem}[Freezing lemma]
Sia $\phi(X,Y)$ una funzione di due variabili. Il valore atteso condizionato di $\phi(X,Y)$ data $Y$ si ottiene "congelando" $Y$ al suo valore osservato e calcolando la media solo rispetto a $X$:
\[ \E[\phi(X,Y)|Y](\omega) = \E[\phi(X, y)]\Big|_{y=Y(\omega)} \]
\end{theorem}
